\documentclass[uplatex,a4paper,11pt]{jsarticle}
\usepackage[T1]{fontenc}
\usepackage{lmodern}
\usepackage{amsmath,amssymb}
\usepackage{longtable,array,booktabs}
\usepackage[margin=23mm]{geometry}
\usepackage{listings}
\usepackage[dvipdfmx,hidelinks]{hyperref}
\usepackage{pxjahyper}
\newcommand{\code}[1]{\texttt{\detokenize{#1}}}
\newcommand{\src}[2]{\par\noindent{\small 参照：\code{#1}、#2行付近。}\par}
\lstset{basicstyle=\small\ttfamily,columns=fullflexible,keepspaces=true,
  breaklines=true,frame=single,numbers=none}
\setlength{\emergencystretch}{3em}
\title{SALMON2 擬ポテンシャル読み込み・前処理の解析\\
  \large \code{s_pp_info} のメンバ、動径範囲、入力形式による相違}
\author{ソースコード調査レポート}
\date{2026年9月24日}
\begin{document}
\maketitle
\begin{abstract}
本稿は、\code{src/common/structures.f90} の \code{s_pp_info} と、
\code{src/atom/pp/input_pp.f90}、\code{src/atom/pp/prep_pp.f90} を中心に、
擬ポテンシャル情報の読み込み、動径テーブル生成、実空間格子への展開を整理する。
とくに、配列容量を表す \code{lmax,lmax0,nrmax,nrmax0}、入力データ範囲を表す
\code{mr}、非局所部分の有効範囲を表す \code{nrps} を区別する。
また、ファイル内に波動関数・密度が存在することと、読み込み実装がそれを保存・利用することを区別し、
PSP8 の価電子密度が保存されない点などを記録する。
\end{abstract}
\tableofcontents
\clearpage

\section{調査範囲と読み方}
対象は調査時点の作業ディレクトリ内のソースである。中心となる3ファイルに加え、初期化元の
\code{salmon_pp.f90}、形式別の \code{read_ps_upf_module.f90} と
\code{read_paw_upf_module.f90}、密度生成の \code{prep_density_pp.f90}、
\code{src/gs/initialization_dft.f90} と \code{src/gs/init_gs.f90} を確認した。
本稿は静的なデータフロー解析であり、各形式の入力ファイルを用いた数値実行試験ではない。
行番号は調査時点の目安であり、変更後にはルーチン名で照合する必要がある。

以下では、\code{ik} を元素種、\code{ia} を原子、\code{i} を動径添字、
$\ell$ を角運動量、$p$ を射影子の通し番号とする。
ソースでは角運動量に \code{ll}、射影子番号にも \code{l} が使われる箇所がある。
\code{nproj(ell,ik)} は磁気量子数による縮退を含まない動径射影子数である。

\section{処理全体の役割分担}
\begin{enumerate}
\item \code{salmon_pp.f90: read_pslfile -> init_pp}：
配列確保、元素記号・質量、容量、初期値を設定する。一部形式ではヘッダを事前に読む。
\item \code{input_pp.f90: input_pp}：ルートMPIプロセス上で元素種ごとに
形式別読み込み、カットオフ設定、規格化、maskingの有無に応じたテーブル生成を行う。
終了時に主要な元素別テーブルをMPI broadcastする。
\item \code{prep_pp.f90: init_ps}：動径テーブルと原子位置を用い、格子点の列挙、
非局所射影子の補間、局所ポテンシャルの生成を行う。
\end{enumerate}

\code{prep_pp.f90} の \code{pp} は \code{intent(in)} である。
このファイルは \code{s_pp_info} を更新せず、結果を \code{s_pp_grid} 型の
\code{ppg}、\code{Vpsl}、必要な作業領域へ書き込む。
したがって「prep段階でppのメンバに情報が追加される」という理解は正しくない。
\src{src/atom/pp/salmon_pp.f90}{22}
\src{src/atom/pp/input_pp.f90}{28}
\src{src/atom/pp/prep_pp.f90}{22}

\section{容量・入力範囲・有効範囲の区別}
\subsection{6変数の意味と決定元}
\begin{longtable}{@{}p{25mm}p{32mm}p{98mm}@{}}
\toprule 変数 & 決定元・現在の値 & 意味\\\midrule\endhead
\code{lmax0} & \code{init_pp} 内の定数、8 & 読み込み作業配列の角運動量容量。入力の最大角運動量の検査にも使う。\\
\code{lmax} & 呼び出し元から設定、8 & 元素別テーブルと格子展開用配列の角運動量容量。実際の元素別上限は \code{mlps(ik)}。\\
\code{nrmax0} & \code{init_pp} 内の定数、50000 & \code{upp,vpp} などの動径作業配列の添字上限。動径添字は \code{0:nrmax0}。\\
\code{nrmax} & 孤立系20000、周期系3000 & \code{rad} と元素別動径テーブルの確保長。ファイル外側への動径座標延長の範囲にも使う。\\
\code{mr(ik)} & ファイルの点数・読み込み件数から決定 & 入力由来の動径データ範囲。形式により点数との対応が異なり、常に点数または最終添字のどちらか一方とはいえない。\\
\code{nrps(ik)} & 読み込み後、\code{rps} と \code{rad} から計算 & 非局所射影子テーブルの有効範囲の終端。masking時には再計算する。\\
\bottomrule
\end{longtable}

\subsection{処理段階との対応}
\begin{longtable}{@{}p{24mm}p{42mm}p{43mm}p{46mm}@{}}
\toprule 変数 & 読み込み & 前処理・生成 & 実空間格子への展開\\\midrule\endhead
\code{lmax0} & 作業配列容量と入力検査 & 一時配列容量。値は不変 & \code{prep_pp} で直接使用しない\\
\code{lmax} & 保存容量と入力検査 & 射影子テーブル容量。値は不変 & \code{ppg} の配列容量\\
\code{nrmax0} & 作業配列容量と入力検査 & masking用一時配列など。値は不変 & \code{prep_pp} で直接使用しない\\
\code{nrmax} & 動径座標・保存テーブル容量 & 外側への延長、\code{nrps} の探索・検査 & 補間係数の確保長、孤立系局所項の探索範囲\\
\code{mr} & 読み込みループ範囲 & 規格化、微分、密度・テーブル生成範囲 & \code{prep_pp} では直接使わない。初期密度生成では積分範囲\\
\code{nrps} & 通常は直接読まない & カットオフから生成・maskingで更新 & 非局所補間係数の作成と区間探索\\
\bottomrule
\end{longtable}

これらは3次元実空間格子の点数ではない。実空間格子は \code{lg,mg} が管理する。
原子ごとの非局所領域に入る実空間格子点数は \code{ppg}\texttt{\%}\code{mps(ia)}、その容量は
\code{ppg}\texttt{\%}\code{nps} である。

現在は \code{lmax=lmax0=8} だが意味は異なる。
射影子形式では作業配列の第2添字は角運動量そのものではなく射影子番号となり、
最大2個の動径射影子を各角運動量に想定して、おおむね $2(\code{lmax0}+1)$ 個分を確保する。
局所ポテンシャル用にさらに1列を持つ \code{vpp} の確保範囲は
\code{(0:nrmax0,0:2*lmax0+2)} である。

\subsection{形式によるmrの違い}
\begin{longtable}{@{}p{35mm}p{120mm}@{}}
\toprule 形式 & 本読み込み後の設定\\\midrule\endhead
KY & ヘッダから \code{mr} を読み、\code{0:mr} の値を読む。\\
ABINIT / PSP8 & \code{mr=mmax-1}。PSP8は \code{mmax} 個の動径点を読む。\\
ABINITFHI / FHI & \code{mr=mr_l}。正半径のデータに加え、原点の波動関数・ポテンシャルをプログラム側で設定する。\\
ADPACK & 読み込み件数 \code{icount} を設定し、原点を追加する。\\
通常UPF & 入力の最初の半径が0なら \code{mr=nrr}。そうでなければ原点を追加して \code{mr=nrr+1}。\\
\bottomrule
\end{longtable}

事前ヘッダ読み込みと本読み込みで値が変わる場合もある。例えばABINIT系の事前読み込みは
\code{mmax} を入れるが、本読み込みでは \code{mmax-1} になる。
各ループの上下限と配列の添字規約を合わせて読む必要がある。

\subsection{添字、外側補完、nrpsの生成}
基本的な座標対応は次のとおりである。
\begin{lstlisting}
upp(i,p), vpp(i,p)       <-> rad(i+1,ik)
udvtbl(i,p,ik)           <-> radnl(i,ik)
vloctbl(i,ik)            <-> rad(i,ik)
\end{lstlisting}
KYは等間隔格子、ABINITは形式固有の式、FHI系は比率による格子延長などを用いる。
読み込み後の共通処理には、\code{mr+1:nrmax} について
\code{upp} を0、\code{vpp} の \code{0:mlps} 列を \code{-zps/rad(i)} とする処理がある。
ただし、射影子形式ではこれらの列がポテンシャルでない場合があり、局所列も末尾に移動する。
この共通処理を「全形式の局所ポテンシャルを正しく延長する処理」と一律には解釈できない。

チャネル別カットオフ \code{rrc} が非ゼロなら、その最大値を \code{rps} に設定する。
その後、\code{rad(i)>rps} を初めて満たす \code{i} を \code{nrps} とする。
\code{nrps>=nrmax} は停止条件である。
\code{nrloc=nrps}、\code{rloc=rps}、\code{radnl=rad} もこの段階で設定する。
\src{src/atom/pp/input_pp.f90}{103}

maskingでは \code{rps=gamma_mask*rps} として \code{nrps} を再探索し、
\code{rps=rad(nrps)} と格子点に合わせる。\code{mr,nrmax,nrmax0} や座標配列自体は変えない。
\code{nrloc,rloc} はmasking前の値を保持する。
なお \code{nrmax0} の検査だけで、より短い \code{rad(nrmax,...)} などの
配列境界がすべて保証されるわけではない。
\src{src/atom/pp/input_pp.f90}{1324}

\section{s\_pp\_infoの全メンバの役割と設定箇所}
容量・範囲の6変数は前節で説明した。以下に残りのメンバをまとめる。
\src{src/common/structures.f90}{229}

\subsection{元素・チャネル情報}
\begin{longtable}{@{}p{39mm}p{116mm}@{}}
\toprule メンバ & 役割と書き込み\\\midrule\endhead
\code{zion} & 読み込み中の擬イオン電荷の実数スカラー。ABINIT、PSP8、ADPACK、UPFで設定。元素別保存先ではなく、KY/FHI系は経由しない。\\
\code{zps(:)} & 元素別の整数擬イオン電荷。各形式のヘッダ等から設定し、クーロン項や初期密度の電子数に使う。\\
\code{atom_symbol(:)} & \code{init_pp} が \code{izatom} に応じて設定。ログとファイル名に使用。\\
\code{rmass(:)} & \code{init_pp} の元素表で設定。読み込み後に \code{system}\texttt{\%}\code{mass} へコピー。\\
\code{mlps(:)} & 各元素で使う最大角運動量。ファイルから設定。KY、ABINIT、ABINITFHI、FHIでは非負の \code{Lmax_ps} 指定で上書きする。\\
\code{lref(:)} & 局所ポテンシャルの列番号。初期値は \code{lloc_ps}。PSP8、ADPACK、UPFでは局所項を独立した末尾列に置き、その列番号へ変更する。\\
\code{nproj(:,:)} & 角運動量ごとの動径射影子数。初期値0。射影子形式で読み込み・集計し、全列0なら \code{0:mlps} を各1に設定。\\
\code{num_orb(:)} & 初期値0。PAW-UPF事前読み込みで軌道／射影子チャネル数を設定し、PAW側の配列確保に利用。\\
\bottomrule
\end{longtable}

\subsection{座標と局所・非局所の範囲}
\begin{longtable}{@{}p{39mm}p{116mm}@{}}
\toprule メンバ & 役割と書き込み\\\midrule\endhead
\code{rad(:,:)} & 元素別動径座標。形式別に読み込み・生成・延長。\\
\code{radnl(:,:)} & 非局所補間用動径座標。共通処理で \code{rad} をコピー。独立した再メッシュはしない。\\
\code{rps(:)} & 非局所射影子の空間範囲。通常は \code{rrc} の最大値。PAW経路では直接設定も行う。maskingで更新。\\
\code{rloc(:)}、\code{nrloc(:)} & masking前の \code{rps,nrps} を保存。局所項のフーリエ変換は \code{nrloc} を積分終端とする。\\
\code{rps_ao(:)}、\code{nrps_ao(:)} & DFT+U用原子軌道範囲。masking前の \code{rps,nrps} をコピー。\code{upp} の非ゼロ領域を走査する処理はあるが、走査結果を代入する行はコメントアウトされている。\\
\bottomrule
\end{longtable}

\subsection{作業配列と保存配列}
\begin{longtable}{@{}p{39mm}p{116mm}@{}}
\toprule メンバ & 役割と書き込み\\\midrule\endhead
\code{upp(:,:)} & 元素ごとに再利用する動径波動関数。KY、ABINIT、FHI系で読み込む。maskingなしでは半径・角度因子を掛けて上書き。射影子形式は原則これを読まない。\\
\code{vpp(:,:)} & 波動関数入力形式では角運動量別ポテンシャル。射影子形式では射影子と独立した局所列を格納。\\
\code{dupp(:,:)} & maskingなしで、変換後の \code{upp} の動径微分を差分計算。\\
\code{dvpp(:,:)} & maskingなしで \code{vpp} の動径微分を差分計算。形式により微分対象がポテンシャルまたは射影子となる。\\
\code{upp_f(:,:,:)} & 各元素の処理末尾で \code{upp} をコピー。maskingなしでは変換済み、ありでは \code{upp} を直接変換しないため、内容の意味が経路で異なる。\\
\code{vpp_f(:,:,:)} & 各元素の処理末尾で \code{vpp} をコピー。孤立系局所項の補間で \code{lref} 列を使用。\\
\code{vpp_so(:,:)} & ADPACK、SO対応UPFから読むSO用追加射影子。\\
\code{dvpp_so(:,:)} & maskingなしのSO処理で \code{vpp_so} の微分を設定。\\
\code{vpp_f_so(:,:,:)} & 宣言のみ。調査した \code{src} 内に確保・代入・利用が見当たらない。\\
\bottomrule
\end{longtable}
maskingありでは微分をローカル作業配列に計算するため、\code{dupp,dvpp} に
意味のある結果が常に存在すると考えてはいけない。
\src{src/atom/pp/input_pp.f90}{234、848、1111}

\subsection{規格化と最終テーブル}
\begin{longtable}{@{}p{39mm}p{116mm}@{}}
\toprule メンバ & 役割と書き込み\\\midrule\endhead
\code{anorm(:,:)} & 共通処理で規格化係数の絶対値の平方根を設定。射影子形式では読み込んだ係数を変換。\\
\code{inorm(:,:)} & 係数の符号 $+1/-1$。元の絶対値が $10^{-10}$ 未満なら0。\\
\code{anorm_so(:,:)}、\code{inorm_so(:,:)} & SO用の対応する係数と符号。SO検出時に変換。\\
\code{vloctbl(:,:)} & \code{making_ps_*} で局所列から生成する局所ポテンシャル。\\
\code{dvloctbl(:,:)} & 同じ段階で計算する局所ポテンシャルの動径微分。\\
\code{udvtbl(:,:,:)} & 規格化と半径・角度因子の処理を済ませた非局所射影子の動径テーブル。\\
\code{dudvtbl(:,:,:)} & 上記の動径微分。\\
\code{udvtbl_so(:,:,:)}、\code{dudvtbl_so(:,:,:)} & \code{making_ps_*_so} が作るSO用の対応するテーブル。\\
\bottomrule
\end{longtable}

読み込み時の波動関数を $u_\ell(r)$、$\Delta V_\ell=V_\ell-V_{\rm loc}$ とすると、
ポテンシャル・波動関数入力では、コードの離散和により
\begin{equation}
 A_\ell\simeq\int u_\ell(r)^2\Delta V_\ell(r)\,dr,
 \qquad a_\ell=\sqrt{|A_\ell|},\qquad s_\ell=\operatorname{sgn}(A_\ell)
\end{equation}
を計算し、\code{anorm} に $a_\ell$、\code{inorm} に $s_\ell$ を入れる。
ゼロ判定された項は後段で除外される。maskingなしの非ゼロ項のテーブルは
\begin{equation}
 T_\ell(r)=\frac{u_\ell(r)\Delta V_\ell(r)}{a_\ell}
 \frac{\sqrt{(2\ell+1)/(4\pi)}}{r^{\ell+1}}
\end{equation}
に相当する。原点では別途コピー・外挿処理を行う。

射影子入力では、読み込んだ係数 $D_p$ に対し $a_p=\sqrt{|D_p|}$ とする。
形式別変換後の \code{vpp} 内の射影子を $b_p(r)$ と書けば、maskingなしでは
\begin{equation}
 T_p(r)=a_p b_p(r)\frac{\sqrt{(2\ell+1)/(4\pi)}}{r^{\ell+1}}
\end{equation}
となる。前者は規格化係数で割り、後者は掛ける点が異なる。
maskingありでは $u\Delta V$ または射影子を一時配列に取り、フーリエ空間の
フィルタリングとマスク処理を行ってからテーブル化する。
\src{src/atom/pp/input_pp.f90}{147、848、1200}

\subsection{密度・NLCC・原子軌道}
\begin{longtable}{@{}p{39mm}p{116mm}@{}}
\toprule メンバ & 役割と書き込み\\\midrule\endhead
\code{rho_pp_tbl(:,:)} & 初期電子密度用の動径分布。$dr$積分が電子数に対応する量で、3次元密度そのものではない。通常UPFから読むか、条件付きで \code{upp} から生成する。\\
\code{flag_nlcc} & 冒頭で偽にし、各元素のNLCCフラグとの論理和で更新する。\\
\code{rho_nlcc_tbl(:,:)} & \code{making_ps_*} で0に初期化後、NLCCありなら \code{rhor_nlcc(i-1,0)} をコピー。中心値との比が $10^{-7}$ 未満で終了。\\
\code{tau_nlcc_tbl(:,:)} & 同じループで $\frac14(\rho'_c)^2/\rho_c$ を設定。微分情報が読み込み側で提供されるかに依存する。\\
\code{upptbl_ao(:,:,:)} & maskingなしで変換後の \code{upp(i-1,p)} をコピーするDFT+U用テーブル。maskingありでは対応する代入がない。\\
\code{dupptbl_ao(:,:,:)} & 原子軌道微分用と読める名称だが、\code{init_pp} の確保・ゼロ初期化以外に計算・利用が見当たらない。\\
\bottomrule
\end{longtable}

\section{ファイル形式と情報の読み込み・保存}
\subsection{形式別一覧}
この表は形式の一般仕様ではなく、現在の読み込み実装が行う処理を示す。
\begin{longtable}{@{}p{25mm}p{40mm}p{47mm}p{43mm}@{}}
\toprule 形式 & 波動関数 & 価電子密度 & コア密度\\\midrule\endhead
KY & \code{upp} に読む & \code{rhopp} に読むが永続保存しない & 部分コアのヘッダ情報は読むが補正は未実装\\
ABINIT & \code{upp} に読む & 密度テーブルの読み込みなし & NLCCテーブルを構築しない\\
ABINITFHI & \code{upp} に読む & 読み込みなし & 末尾にあれば密度と微分を読む\\
FHI & \code{upp} に読む & 読み込みなし & このルーチンでは読まない\\
PSP8 & 読まない & 読むが積分値表示のみ。\code{rho_pp_tbl} に保存しない & \code{fchrg>0} なら密度と微分を読む\\
ADPACK & 読まない & 読み込みなし & \code{density.PCC} があれば密度を読む。微分は読まない\\
通常UPF & \code{PP_PSWFC/PP_CHI} を読む処理なし & \code{PP_RHOATOM} を \code{rho_pp_tbl} に保存 & \code{PP_NLCC} を読む。微分の読み込み・生成なし\\
PAW-UPF & \code{PP_CHI} を \code{upp_f}、部分波を \code{paw} 配列へ & \code{paw}\texttt{\%}\code{rhops} に読む & \code{paw}\texttt{\%}\code{rhocps/rhocae} に読む\\
\bottomrule
\end{longtable}

KY、ABINIT、FHI系は角運動量別ポテンシャルと波動関数から非局所射影子を構成する。
PSP8、ADPACK、UPFは射影子・係数・局所ポテンシャルを入力として使用する。
例えばADPACKは射影子に半径を掛け、UPFはエネルギー単位の変換を含む処理を行うため、
\code{vpp} をファイルの生データと同一視することもできない。
\src{src/atom/pp/input_pp.f90}{288、343、402、498、614、686}
\src{src/atom/pp/read_ps_upf_module.f90}{274、297}

\subsection{PSP8の密度は読み捨てられる}
PSP8の価電子密度部分は、次のように局所変数へ読み、積分値だけを表示する。
\begin{lstlisting}[language=Fortran]
sum_rho_pp=0.0d0
do i=1,pp%mr(ik)+1
  read(4,*) dummy_text, r_tmp, rho_tmp
  sum_rho_pp=sum_rho_pp+rho_tmp*r_tmp**2
end do
write(*,*) "sum(rho_pp)=",sum_rho_pp*dr
\end{lstlisting}
\code{rho_pp_tbl} への代入は存在しない。したがって、ファイルに密度があるにもかかわらず、
本実装ではその情報を初期密度構築に使えない。波動関数の読み込みも行わない。
\src{src/atom/pp/input_pp.f90}{577}

\subsection{波動関数からの密度生成}
ポテンシャル・波動関数入力で、maskingなし、かつ \code{method_init_density/='wf'}
の場合、次の蓄積を行う。
\begin{equation}
 \code{rho_pp_tbl}(i,ik)\mathrel{+}=(2\ell+1)u_\ell(i)^2.
\end{equation}
蓄積した電子数が \code{zps} を超えると角運動量ループを終了する。
これはファイルの占有数を忠実に読み込む処理ではなく、実装固有の蓄積・打ち切りである。
maskingありには対応する密度生成がない。通常UPFは読み込み時に密度を設定するので、
この波動関数からの生成には依存しない。
\src{src/atom/pp/input_pp.f90}{1123}

\subsection{初期密度方式と、波動関数がない場合}
射影子入力では、UPF以外で \code{method_init_density} が \code{wf} または
\code{read_dns_cube} でなければ、\code{input_pp} が次のメッセージで停止する。
\begin{lstlisting}
radial density is not available (method_init_density=pp...)
\end{lstlisting}
PSP8の場合、この制限はファイルに密度が存在しないからではなく、読み込み実装が
利用可能な密度テーブルを保存していないことに対応する。
\src{src/atom/pp/input_pp.f90}{92}

\code{method_init_density='wf'} は、擬ポテンシャルファイルの \code{upp} を意味しない。
Gaussianや乱数等で初期化した計算対象の軌道 \code{spsi} から密度を求める。
したがって、原子波動関数を含まないPSP8でもこの方式は使える。
\code{read_dns_cube} は別ファイルから密度を読み、\code{pp} / \code{pp_magdir} は
\code{rho_pp_tbl} に基づいて密度を作る。
\src{src/gs/initialization_dft.f90}{266}
\src{src/gs/init_gs.f90}{40}

\code{calc_density_pp} は \code{1:mr(ik)} を積分範囲として動径密度をフーリエ変換し、
原子位置の位相を掛けて逆変換する。密度の積分範囲は \code{nrps} ではない。
$G=0$ 成分は \code{zps} から直接設定するため、密度テーブルが0でも
電子数成分だけは非ゼロとなり得る。
\src{src/atom/pp/prep_density_pp.f90}{40}

\section{prep\_ppにおける利用先}
\begin{longtable}{@{}p{39mm}p{116mm}@{}}
\toprule 処理 & 読み取りと出力\\\midrule\endhead
\code{calc_nps/calc_jxyz} & \code{rps} 内の実空間格子点を列挙し、\code{ppg}\texttt{\%}\code{mps,jxyz,rxyz} などを構築。\\
\code{set_lma} & \code{mlps,nproj,inorm} を参照し、無効項を除いて原子・射影子・磁気量子数に通し番号を付ける。出力は \code{ppg}\texttt{\%}\code{nlma,lma_tbl,ia_tbl}。\\
\code{calc_uv} & \code{radnl,udvtbl} の \code{1:nrps} から初回にスプライン係数を作る。各格子点で補間し、\code{ylm(x,y,z,ll,m)} を掛けて \code{ppg}\texttt{\%}\code{uv} を構築。\code{ppg}\texttt{\%}\code{rinv_uvu=inorm*hvol}。\\
孤立系局所項 & \code{rad} と \code{vpp_f(:,lref,ik)} を線形補間し、\code{Vpsl}\texttt{\%}\code{f} と \code{ppg}\texttt{\%}\code{Vpsl_ion} を生成。\\
周期系局所項 & \code{rad,vloctbl,zps} と終端 \code{nrloc} で $V_{\rm loc}+Z/r$ の短距離部分を動径積分。クーロン項と原子位相を加え、逆変換して \code{Vpsl}\texttt{\%}\code{f} を生成。\\
SO / DFT+U & 別モジュールの \code{calc_uv_so/calc_uv_plusU} に \code{pp} を渡し、\code{ppg} 側のデータを生成。\\
\bottomrule
\end{longtable}
\code{dudvtbl,dvloctbl} は構築されるが、このファイルの通常の格子展開では直接参照しない。
非局所スプライン係数は \code{udvtbl} の値から生成する。
\src{src/atom/pp/prep_pp.f90}{372、423、555、802}

\section{実装上の注意点と確認すべき事項}
以下は代入・分岐の追跡から確認した事項である。特定入力での数値的影響は別途実行検証を要する。
\begin{enumerate}
\item \textbf{作業配列と形式フラグの持ち越し。}
\code{upp,vpp,vpp_so} は初期確保時にゼロ化されるが、元素ループ冒頭で一律リセットされない。
\code{flag_beta_proj_is_given,flag_so} もモジュール変数で元素ごとにリセットされない。
形式を混在させると前元素の値や分岐状態が影響し得る。「読まない情報は常に0」とはいえない。

\item \textbf{波動関数の存在とAOテーブル。}
通常UPFは \code{PP_CHI} を読まず、射影子から原子波動関数を復元する処理もない。
maskingなしでは \code{upp} からAOテーブルへのコピーが実行されても、有効な原子波動関数が
存在する保証にはならない。maskingありにはAOテーブル生成自体がない。

\item \textbf{PAW-UPFの保存経路の不整合。}
専用リーダは \code{PP_CHI} を \code{upp_f} に直接読むが、呼び出し元の末尾が
\code{upp_f(:,:,ik)=upp(:,:)} と上書きする。
また、価電子密度とコア密度は \code{paw} 側に保存され、通常の
\code{rho_pp_tbl} や \code{rhor_nlcc} へのコピーはこの経路にはない。
通常UPFとPAW-UPFを同じ密度利用経路として扱うことはできない。

\item \textbf{密度データ存在の検査。}
通常UPFの読み込み用 \code{cdd} はゼロ初期化される。
\code{PP_RHOATOM} が欠けるとゼロのままであり、初期密度方式の許可は
実データの存在ではなく \code{ps_format='UPF'} で判定される。

\item \textbf{NLCC密度微分。}
\code{tau_nlcc_tbl} は \code{rhor_nlcc(:,1)} に依存する。
ADPACKと通常UPFは密度微分を読まず、共通一時配列は元素ごとに一律クリアされない。
そのため、密度のみの入力から運動エネルギー密度が適切に構築されるとは限らない。

\item \textbf{SOとmaskingによるカットオフ更新。}
通常成分とSO成分の両方が \code{ps_masking} を呼ぶ。
それぞれが同じ \code{pp}\texttt{\%}\code{rps} を更新するため、SOとmaskingがともに有効な経路では
カットオフの拡大・格子点への丸めが二度行われる。

\item \textbf{maskingの表示と実処理。}
ログには \code{anorm,inorm} が変更される旨の文言があるが、maskingルーチンには
両者を再計算する代入がない。出力メッセージだけから再規格化を推測してはいけない。

\item \textbf{未使用メンバと有効範囲。}
\code{vpp_f_so} は宣言のみ、\code{dupptbl_ao} は確保・ゼロ化のみである。
また、配列の確保長と意味のある値が書き込まれた範囲は一致しない。
\code{mr,nrps,nrloc} と形式別のループ範囲を合わせて判断する必要がある。
\end{enumerate}
\src{src/atom/pp/read_paw_upf_module.f90}{140}
\src{src/atom/pp/input_pp.f90}{20、234、1018、1302}

\section{再読時の要点}
\begin{itemize}
\item 容量は \code{lmax/lmax0/nrmax/nrmax0}、入力由来の範囲は \code{mr}、
非局所有効範囲は \code{nrps}、実空間格子点数は \code{ppg}\texttt{\%}\code{mps/nps} と区別する。
\item \code{input_pp} が元素別動径情報を構築し、\code{prep_pp} はそれを読む。
\item \code{vpp} は形式によってポテンシャルにも射影子にもなる。
\item ファイルに含まれる情報、リーダが読む情報、保存される情報、後段が利用する情報を分ける。
PSP8の価電子密度はこの区別が必要な代表例である。
\item 初期密度方式 \code{wf} は計算対象の軌道からの密度生成であり、原子波動関数 \code{upp} の使用ではない。
\end{itemize}

\appendix
\section{コンパイル方法}
up\LaTeX{} とdvipdfmxを用いる。目次・相互参照のため2回コンパイルする。
\begin{lstlisting}
uplatex -interaction=nonstopmode -halt-on-error report_pp.tex
uplatex -interaction=nonstopmode -halt-on-error report_pp.tex
dvipdfmx report_pp.dvi
\end{lstlisting}
\end{document}
